\section{Comparing the good spaces on different real forms}
\label{sec:trace}

Let $V_s,V_{s'}$ be two of the rank-$r$ hermitian spaces $Y_0,W_j$
with the prescribed finite frames. These frames identify their groups of
finite adelic points and their principal level $K$. They have a common
doubled Harish--Chandra parameter $l$ and possibly different real
signatures $s,s'$. Write $M_s$ for the space of forms of type $l$ at
level $K$, $e_0$ for the good projector and $M_s^g=e_0M_s$. The Hecke
algebra $\mathcal H$ used here consists of the spherical operators away
from the chosen places, acting simultaneously on the different spaces.
For its elements the desired equality is
\begin{equation}\label{tr:good}
  \operatorname{tr}(T\mid M_s^g)=\operatorname{tr}(T\mid M_{s'}^g)
  \qquad(T\in\mathcal H).
\end{equation}
We derive this under the geometric inputs specified below. The spreading
argument needs $T=1$ at $j=m$. The chain determinant also needs a common
Hecke operator at each $j$.

\subsection{Getting a cuspidal component from a label}\label{tr:label}
At the chosen split places the principal level is the kernel of reduction
onto the product of the finite groups $G_w=\mathrm{GL}_r(k_w)$.
The good idempotent is $e_0=\sum_c e_c$, where $c$ ranges over the good
tuples and $e_c$ is the tensor product of their character idempotents.
Each good tuple has a primitive degree-$r$ label at some $w$ above
the first level prime. Its root in $\mathbb F_{q^r}$,
$q=|k_w|$, generates $\mathbb F_{q^r}^{\times}$. Under the chosen
identification with its character group this gives a character of order
$q^r-1$. Its Frobenius orbit has length $r$, since $0<q^h-1<q^r-1$
for $0<h<r$. The irreducible factor of the label therefore indexes
Green's cuspidal character $\chi$ of $G_w$, corresponding to the degree-$r$
orbit and the one-box partition~\cite[\S8, Theorem~14]{green}.

One can also read off the invariant multiplicity from the polynomial
definition of the label. In the Hall algebra let $*$ be parabolic
induction, $\mathbf{u}$ the constant function, $\mu$ its $*$-inverse,
and $\beta$ the sum of the chosen multiplicative character over
the eigenvalues with multiplicity. Green's function for the irreducible
label is the rank-$r$ part of $C=(\beta\mu)*\mathbf{u}$ up to sign,
with $\beta\mu$ the pointwise product. Block averaging $\delta$ is
compatible with $*$ by the finite Mackey formula (partition a stable
subspace by its intersection with the first block and its image in the
quotient). Since $\delta\mathbf{u}=\mathbf{u}\otimes\mathbf{u}$,
one has $\delta\mu=\mu\otimes\mu$. The eigenvalue sum is constant on
a block extension and equals the sum on its diagonal blocks. Consequently
\[
 \delta(\beta\mu)=(\beta\mu)\otimes\mu+\mu\otimes(\beta\mu),
 \qquad \delta C=C\otimes 1+1\otimes C,
\]
where $1$ is supported in rank zero. Taking two identity diagonal
blocks of positive ranks gives zero average on their radical. This
average of the character is the dimension of the invariant subspace.

Here is the vanishing that the trace comparison requires. Let $N$ be
the unipotent radical of a proper coordinate maximal parabolic, and
let $\pi$ afford $\chi$. The averaging operator
$|N|^{-1}\sum_{u\in N}\pi(u)$ projects onto $\pi^N=0$ by cuspidality.
Multiplication on the left and right and taking the trace gives
$\sum_{u\in N}\chi(aub)=0$ for every $a,b\in G_w$. The coefficients of
the central idempotent are $\chi(1)\overline{\chi(g)}/|G_w|$.
Their sums over these translates also vanish. The map
$Z\mapsto 1+pZ(1-p)$, with $p$ the projection onto the chosen
coordinate subspace, covers $N$ with constant fibre size, so summing
over all matrices $Z$ gives the same vanishing. In $e_c$ the remaining
finite factors are constant during this summation.

Lift the tuple idempotent by a section of reduction and convolve with
$\mathbf 1_K$. On the standard maximal compact this function is the
coefficient function of $e_c$ evaluated at the reduction: precisely the
chosen lift of that reduction contributes. Outside the compact it
vanishes. Right multiplication by the local integral lifts of $N$ stays
in the same compact coset. Thus all the local radical sums that define
cuspidality at $w$ vanish. Convolution on the left by the group-algebra
representative of $T$ preserves these sums. This proves the local
hypothesis for the finite test function of $T e_c$.

\subsection{Comparing traces with two local conditions}\label{tr:geometry}
For clarity, the geometric branch of the argument uses the following
inputs for the kernels and Haar normalizations of these spaces.
They are hypotheses of the present reconstruction.
\begin{itemize}
\item[(i)] At a common compact open level $K'$ for either signature,
the kernel for $l$ has cusp decay of order $10r^5$, and the convergent
pretrace expansion for a level-preserving test is
$\operatorname{vol}(K')\operatorname{tr}=\sum_R I_s(R)$.
Here $I_s(R)$ is the contribution of the characteristic polynomial
$R\in E[X]$ of rational elements. The expansions unfold to class
orbital integrals for irreducible $R$, with absolutely integrable cyclic
triple-kernel integrals also after weighting by a one-step logarithmic
norm increment.
\item[(ii)] For coprime nonconstant local factors, a conjugacy-invariant
distribution supported on the characteristic-polynomial fibre of their
product kills a test function with vanishing radical sums at the split
place. This is the local parabolic descent step.
\item[(iii)] For irreducible $R$ with all roots on the unit circle at
every infinite place, each elliptic class integral factors into a finite
factor and the product of real orbital integrals. Changing the signs
attached to the eigenlines, with total sign product fixed at $1$, gives
a bijection of the corresponding frame classes preserving the finite factor.
For regular $z=(z_i)$ on the compact torus, the real kernels obey
\[
 (-1)^{n_+n_-}\sum_{\sigma\in S_r} O_s(z\circ\sigma)
 =|S_{n_+}\times S_{n_-}|\,\overline{\chi_{l,\mathrm{cpt}}(z)}.
\]
Here $n_+,n_-$ count the two signs, $O_s$ uses the
formal-degree-normalized discrete-series kernel of $l$, and the right side
uses the lowest $K$-type on the compact form at that place.
\end{itemize}
Section~\ref{tr:parabolic-descent} proves (ii) for the split test just constructed. The remaining assumptions concern the indicated kernels on the adelic quotient,
including the finite centralizer measure in (iii).

Suppose the finite test is cuspidal at $w$ and its support at an auxiliary
place consists of integral matrices having zero as a simple root of the
reduced characteristic polynomial. The latter condition excludes a proper
power $R=Q^h$, $h>1$: a monic $Q$ with $Q^h$ integral has integral
coefficients by integral closedness, and its reduction would give zero
with multiplicity divisible by $h$. Each remaining reducible $R$ in the
support has two coprime nonconstant factors over $E$ obtained by
partitioning its primary factors. They stay coprime over the completion.
The functional on the test defined by $I_s(R)$ is conjugation invariant
by change of variables and supported on this fibre. Hence (ii) kills it.

For irreducible $R$, fix a root $z$ with $|z|\ne1$ at one infinite
place, if there is one. Diagonalize a rational representative $\gamma$
there, take its left eigenrow $v$ for $z$, and put
$h(x)=\log\|vx\|$. Then $h(\gamma x)=h(x)+\ell$ with
$\ell=\log|z|$. The centralizer scales this eigenrow, so differences
of $h$ descend under its action. Write $p$ for the convolution
idempotent kernel at infinity. In the unfolded triple integral for the
orbital integral of $p*p*p$ the three successive increments of $h$ from
$x$ to $\gamma x$ sum to $\ell$. Their integrals agree by cyclic
rotation. The integral of the sum of the first two equals that of the
first alone by collapsing the first two kernels using $p*p=p$.
The second integral is therefore zero. Assumption (i) justifies these
integrations and shows that $\ell I_s(R)=0$ after summing the classes.

In the elliptic case, sum (iii) over all orders of the roots at each
infinite place. For fixed signature the eigenline signs of a given
class occur $|S_{n_+}\times S_{n_-}|$ times at each place.
These multiplicities cancel the same factors in the real orbital formula.
The field has an even number $f$ of infinite places, so the total sign
$(-1)^{fn_+n_-}$ is $1$. Every rearrangement of the signature at all
places has total sign product $1$ for the same reason. Matching with
the all-positive signature in (iii) leaves one common finite sum
$C(R)$, and
\[
 I_s(R)=C(R)\prod_{v\mid\infty}\overline{\chi_{l_v,\mathrm{cpt}}(z_v)}
       =I_{s'}(R).
\]
Summing over $R$ and cancelling the positive common volume proves the
two-place trace comparison.

The hypotheses on the data here are quantitative. The field contains
$\mu_{5040}$ because $5040\mid18480\cdot3\mid q_1$ in its cyclotomic
construction. The tower exponent is positive, which gives the parity of
$f$. The admissible ranks give $r\ge2$. At every infinite place the
entries $l_i$ are distinct and congruent to $r-1\pmod2$. For indices of
opposite sign one has
$|l_i-l_j|\ge2(500r^3+10r^5)$ for both signatures. This follows from
the gap in the admissible rows: at $Y_0$ use the initial grid gap; at
$W_j$ the two groups of undoubled parameters lie on opposite sides of
the cut $-\mathrm{liftWeight}_j$ with the prescribed half-gap on each
side. The order $10r^5$ used in (i) is at least $13r^3$. The finite
frames, principal level, and parameter are the same in both applications.

\subsection{Removing the auxiliary local condition}\label{tr:line}
Fix a group-algebra representative $a$ of $T e_c$. Choose a split
place $v$ away from $w,\bar w$, the finite support of the level, and
all nonintegral components of the finitely many elements occurring in
$a$. There are infinitely many split places. At $v$ cut $K$ to the
inverse image of the stabilizer $J$ of the first coordinate line in
$\mathrm{GL}_r(\mathbb F_q)$. Write $E_0$ for averaging over the full
local compact and $A$ for the action of the part of $a$ away from $v$
on this finer level. On the range of $E_0$, $A$ is precisely the old
action of $a$ because the removed components are integral.
Let $P,Q$ denote the unnormalized double-coset operators for
$\operatorname{diag}(\varpi,1,\ldots,1)$ and its inverse at level $J$.
They commute with $A$, and $Q=P^*$ since the inverse double cosets
have the same degree $q^{r-1}$.

Here is the local count. The right cosets for $P$ are indexed by
linear forms $\xi$ with $\xi(e_1)=1$, and those for $Q$ by columns
whose first coordinate is $1$. Representatives are the inverse lift of
the identity with first row replaced by $\xi$, times the diagonal matrix,
and the inverse diagonal matrix times a lift of the identity with first
column replaced by the column, respectively. Multiplying cancels the two
diagonal matrices. For fixed $\xi$ the resulting column $z$ ranges
once over $\xi(z)=1$, since the first row of the row matrix is $\xi$.
The resulting translate of a $J$-fixed vector $u$ depends only on the
line of $z$; write it $u_z$. For $z=e_1$ the count is $q^{r-1}$,
for other multiples of $e_1$ it is zero, and for $z$ off the first
line it is $q^{r-2}$ by two independent affine equations. Since
$\sum_{z\ne0}u_z=(q^r-1)E_0u$, we obtain
\begin{equation}\label{tr:parahoric}
 PP^*u=q^{r-2}\{u+(q^r-1)E_0u\}.
\end{equation}
The positive powers of $P$ are supported on matrices reducing to block
form $\left(\begin{smallmatrix}0&*\\0&B\end{smallmatrix}\right)$ with
$B$ invertible. Products and right multiplication by $J$ preserve that
shape, so zero is a simple root. Removing the place $v$ from $a$ and
cutting there preserve the radical sums at $w$: locally at $v$ all
translates in $a$ are integral and at $w$ the last factor has the
vanishing proved above. Equivalently the removed test is constant along
the $v$ component before cutting, and the cut imposes the independent
$J$ indicator there. Left convolution preserves the vanishing for every
positive power. The two-place result therefore equates
$\operatorname{tr}(A^hP^j)$ on the two fine-level spaces for all $h,j\ge1$.

To recover $\operatorname{tr}(AE_0)$, fix $j$ and decompose into generalized
eigenspaces of $A$. On the $\alpha$-space the trace of $A^hP^j$ equals
$\alpha^h\operatorname{tr} P^j$, since $P$ commutes with the nilpotent
part. For the union of nonzero eigenvalues of the two actions, the
matrix $(\alpha^h)_{1\le h\le N}$ is invertible by Vandermonde, with
$N$ the number of such eigenvalues. Thus the positive power sums of $P$
agree on each nonzero $\alpha$-space. Equation~\eqref{tr:parahoric}
makes $P$ invertible there. Positive power sums determine its multiset
of nonzero eigenvalues with multiplicities, by the same Vandermonde
argument. On this space $B_\alpha$ take determinants in
\eqref{tr:parahoric} to get
\[
  |\det(P\mid B_\alpha)|^2
  =(q^{r-2})^{\dim B_\alpha}(q^r)^{\operatorname{rank}(E_0\mid B_\alpha)} .
\]
Both $P$ and $P^*$ preserve $B_\alpha$, so the restriction still has
the stated adjoint. The common dimension and determinant fix the rank.
Summing $\alpha\operatorname{rank}(E_0\mid B_\alpha)$ gives the desired
old-level trace; $\alpha=0$ contributes zero. This removes the
second-place hypothesis for $T e_c$. Finally
$\operatorname{tr}(T\mid M_s^g)=\operatorname{tr}(Te_0\mid M_s)$
because $T$ preserves the range of the idempotent. Summation over $c$
gives \eqref{tr:good}.

\subsection{The two uses of the comparison}\label{tr:uses}
Taking $T=1$ gives $\dim M_s^g=\dim M_{s'}^g=d_T$ when one of the
spaces is $Y_0$. Applying \eqref{tr:good} to every power of $T$
and using Newton's identities gives equal characteristic polynomials
and equal determinants on the good spaces. For an individual Hecke
character, take the finite union of the joint eigencharacters on the
two principal spaces. For each other character choose a Hecke element
separating it from the selected one and rescale the difference with
its value at the other character to give values zero and one.
Multiplying the separating factors gives a common projector onto the
chosen eigenspace, using the simultaneous eigenspace decomposition
for the commuting normal spherical operators. Its trace is the good
multiplicity. Only the identity and the determinant comparison are
required in the applications just described.

For the latter, one must still identify the operators on the two
spaces with the action of a \emph{single} Hecke element. Fix $j$ and
the tame character. In the convergent expression for the matrix
coefficients of the Rallis operator on principal forms each coset term
is a common scalar (the finite vacuum coefficient with its residue
cardinality normalization) times $\langle z,\pi(\gamma)y\rangle$,
with $\gamma$ supported off the chosen level places. Let $K^0$ be the
standard compact in the off-level group. Average $k\gamma$ over
$k\in K^0/(K^0\cap\gamma K^0\gamma^{-1})$. Compact openness gives a
finite quotient of cosets. This normalized spherical operator has the
same matrix coefficient as $\gamma$ on both spaces: $z$ and $y$ are
fixed there and the compact action is unitary. Thus each pair of
terms lies in the joint evaluation image of $\mathcal H$ in the two
endomorphism spaces, expressed in matrix-coefficient coordinates. That
image is a complex linear subspace of a finite-dimensional space, hence
closed. Taking sums preserves membership. The common residue cardinality
multiplying the normalized series restores the two Rallis operators.
Their determinants on the good spaces therefore coincide by
\eqref{tr:good}. This step uses the Rallis series with its specified
normalization and convergence. The analytic identity relating that
operator to the adjoint times the lift is a separate hypothesis in the
chain argument.

\subsection{Eliminating a split characteristic-polynomial fibre}
\label{tr:parabolic-descent}
Here is a local proof of assertion (ii) of \S\ref{tr:geometry}.
Let $F$ be the local field at the chosen split place, $K=\mathrm{GL}_r(\mathcal O_F)$,
and let $R=R_1R_2$ have coprime monic factors of positive degrees $a,b$
and nonzero constant terms. Write $G=\mathrm{GL}_r(F)$, $M=\mathrm{GL}_a(F)\times
\mathrm{GL}_b(F)$ and $U$ for the upper block radical, identified with
the additive space of $a$ by $b$ matrices. Suppose the compactly supported
locally constant test $f_0$ has
\[
 \int_U f_0(k m u k^{-1})\,du=0\qquad (k\in K,\ m\in M).
\]
For the inflated character idempotent this follows from the radical sums
of \S\ref{tr:label}. Indeed $mu\in K$ forces both diagonal blocks to be
integral and invertible, and then $u\in U(\mathcal O_F)$.
Reduction integrates each residue class with the same measure.
The same conclusion holds for any fixed left and right translates of the
test if the finite radical sums vanish for arbitrary left and right
multipliers: by $G=K M U= M U K$ the outside factors reduce to compact
multipliers and an affine change of the variable in $U$.

\begin{lemma}
Every conjugation-invariant distribution on $G$ supported on the fibre of
$R$ annihilates $f_0$.
\end{lemma}
\begin{proof}
On this closed fibre, the ordered complementary subspaces
$\ker R_1(g),\ker R_2(g)$ have dimensions $a,b$. The projectors onto them
are fixed polynomials in $g$ by B\'ezout's identity. The map to complementary
pairs is therefore continuous and equivariant; the space of pairs is
$G/M$. Over the coordinate pair the fibre is the space $X$ of diagonal
blocks $(A,B)$ with their indicated characteristic polynomials. A test
vanishing on the closed fibre is annihilated by the distribution, since
its compact open support misses that fibre. Tests on the fibre extend
locally constantly with compact support, by a finite compact open cover.
Thus Frobenius descent for locally compact totally disconnected spaces
applies to the distribution on this fibre: $G$ and $M$ are both
unimodular. See Aizenbud, Gourevitch, Rallis and Schiffmann,
\emph{Multiplicity one theorems}, Theorem 2.2~\cite{agrs}.

Concretely the descended distribution on $X$ acts on the test
\[
 (A,B)\longmapsto \int_K\int_U
 f_0(k u\operatorname{diag}(A,B)u^{-1}k^{-1})\,du\,dk .
\]
This is the quotient integral in descent: in the Iwasawa coordinates
$g=kum$ one disintegrates Haar measure as $dk\,du\,dm$, up to a fixed
positive scalar, and absorbs $m$ by invariance of the distribution
on $X$. The displayed integral is compactly supported on $X$.
Indeed the compact support bounds the diagonal and off-diagonal blocks
after conjugation by $K$; the inverse of the map $Z\mapsto ZB-AZ$ is
uniformly bounded there. Its determinant, up to sign, is the nonzero
constant $\operatorname{Res}(R_1,R_2)$, as seen by triangularizing
$A,B$ separately over a splitting field. The same bounds make the
integral locally constant and justify testing the distribution under
the integral.

For fixed diagonal blocks, the upper-right entry of the conjugate by
$u(Z)$ is $ZB-AZ$. Changing variables from $Z$ to $A^{-1}(ZB-AZ)$
expresses the inner integral as a constant times the radical integral
in the hypothesis, hence it vanishes. This argument also covers
nonsemisimple diagonal blocks.
\end{proof}

